\chapter{Multiplier (fmul)}
\label{ch:fmul}

\section{Introduction}

\sig{fmul} computes $\mathtt{op\_a}\times\mathtt{op\_b}$. Two \sig{unpack}
modules split the operands. The \sig{fmul\_core} module works in two pipeline stages. The
first normalizes a subnormal operand and computes the sign and exponent of the
product; the second multiplies the mantissas. When both operands are subnormal the
product underflows to zero, so normalizing one is enough. \sig{norm} normalizes the product, and
\sig{pack} rounds it and forms the result, as Figure~\ref{fig:fmul} shows.

\begin{figure}[H]
\centering
\begin{tikzpicture}[node distance=7mm]
  \node[ugblock] (u) {2 $\times$ \sig{unpack}};
  \node[ugblock,right=of u] (a) {\sig{fmul\_core}\\2 cycles};
  \node[ugblock,right=of a] (n) {\sig{norm}\\1 cycle};
  \node[ugblock,right=of n] (p) {\sig{pack}\\1 cycle};
  \draw[ugarrow] ($(u.west)-(9mm,0)$) node[above right,font=\footnotesize\ttfamily,xshift=-1mm] {op} -- (u);
  \draw[ugarrow] (u) -- (a); \draw[ugarrow] (a) -- (n); \draw[ugarrow] (n) -- (p);
  \draw[ugarrow] (p) -- ($(p.east)+(9mm,0)$) node[above left,font=\footnotesize\ttfamily,xshift=1mm] {res};
\end{tikzpicture}
\caption{\sig{fmul} Block Diagram}
\label{fig:fmul}
\end{figure}

\section{Features}

\begin{itemize}
\item Multiplication with a latency of 4 clock cycles.
\item One new operation per clock cycle.
\item Subnormal operands and results.
\item The five rounding modes of Table~\ref{tab:rnd_mode}.
\item NaN and infinity handling, selected by \sig{SPECIAL\_CASES}.
\end{itemize}

\section{Parameters}

\sig{fmul} has the parameters of Table~\ref{tab:common_params}. Its \sig{fmul\_core}
datapath also has the parameter \sig{OUT\_REG} (default 0), which adds an output
register stage when it is not 0, and \sig{PRENORM} (default 0), which enables the
operand normalization. \sig{fmul} sets \sig{PRENORM} to 1, so \sig{norm} shifts the
product left by at most one position. \sig{fpu} sets \sig{OUT\_REG} to 1 and keeps
\sig{PRENORM} at 0, because its \sig{norm} also serves the adder
(Chapter~\ref{ch:fpu}).

\section{Interface Signals}

\sig{fmul} has the signals of Table~\ref{tab:common_sigs}.

\Needspace*{23\baselineskip}
\section{Instantiation}

The example below instantiates \sig{fmul} for binary32. Compile it with the files
of Table~\ref{tab:files}.

\begin{vcode}
fmul #(
   .DATA_W       (32),
   .EXP_W        (8),
   .SPECIAL_CASES(1)
) fmul0 (
   .clk     (clk),
   .rst     (rst),
   .valid_i (valid_i),
   .valid_o (valid_o),
   .rnd_mode(rnd_mode),
   .op_a    (op_a),
   .op_b    (op_b),
   .res     (res)
);
\end{vcode}

\section{Operation}

\sig{fmul} follows the protocol of Section~\ref{sec:protocol}, with a latency of
4 clock cycles. The unit samples the operands in the \sig{valid\_i} cycle, so they
can change in the next cycle. A new operation can start in every clock cycle, and
the results come out in order.

Table~\ref{tab:fmul_special} lists the results for special operands when
\sig{SPECIAL\_CASES} is 1. The sign of an infinite or zero result is the
exclusive OR of the operand signs.

\begin{twotable}{\sig{fmul} Special Results}{tab:fmul_special}{Operands}{Result}{7cm}
Either operand is NaN & NaN \\ \hline
Infinity and zero & NaN \\ \hline
Infinity and a nonzero number & Infinity \\
\end{twotable}
